\documentclass[a4paper,11pt,final]{report}
% Pour une impression recto verso, utilisez plutôt ce documentclass :
%\documentclass[a4paper,twoside,11pt,final]{article}

\usepackage[francais]{babel}
\usepackage[utf8]{inputenc}
\usepackage[T1]{fontenc}
\usepackage[pdftex]{graphicx}
\usepackage{setspace}
\usepackage[colorlinks=true,linkcolor=black,urlcolor=blue]{hyperref}
\usepackage[french]{varioref}
\usepackage[top=7em, bottom=7em, left=7em, right=7em]{geometry}
\usepackage{wrapfig}
\usepackage{fancyhdr}
\fancyhf{}
\renewcommand{\footrulewidth}{0.05em}
\lhead{\leftmark}
\cfoot{\thepage}
\lfoot{\reportauthor}
\rfoot{Projet SMA}

\newcommand{\reporttitle}{Conception d'un système multi-agent dans un environnement 3D}     % Titre
\newcommand{\reportauthor}{Louis \textsc{LE CLEC'H}\\
                           Romain \textsc{SAGEAN}} % Auteur
\newcommand{\reportsubject}{} % Sujet
\newcommand{\HRule}{\rule{\linewidth}{0.5mm}}
\newcommand{\siecle}[1]{\textsc{\romannumeral #1}ème}
\setlength{\parskip}{1ex} % Espace entre les paragraphes


\hypersetup{
    pdftitle={\reporttitle},%
    pdfauthor={\reportauthor},%
    pdfsubject={\reportsubject},%
   % pdfkeywords={rapport} {vos} {mots} {clés}
}

\begin{document}
  \begin{titlepage}
\begin{flushleft} 
  \large
  \emph{Auteur :}\\
  \reportauthor\\
\end{flushleft}
\begin{center}
  
  \begin{center}
    \begin{tabular}{c}
      \\
      \includegraphics [width=55mm]{Images/ENSEIRB.png} \\
      \\
      \includegraphics [width=50mm]{Images/ENSC.jpg}\\
      \\
      \\
    \end{tabular}
    
    
      
    \textsc{\Large \reportsubject}\\[0.5cm]
           {\large 20 Decembre 2013}\\
           
           
           \HRule \\[0.4cm]
                  {\huge \bfseries \reporttitle}\\[0.4cm]
                  \HRule \\[1.5cm]
                  
                  \begin{center}
                    
                    
                    
                    
                  
                  
                  %\vfill                  
                  
                  
                  
  \end{center}    
\end{center}
\begin{flushbottom}
\begin{flushright} 
  \large
  \emph{Responsable :} \\
  M.~Pierre Alexandre Favier\\

\end{flushright}
\end{flushbottom}
\end{center}
\end{titlepage}

  %\cleardoublepage % Dans le cas du recto verso, ajoute une page blanche si besoin
  \tableofcontents % Table des matières
  \listoffigures
 % \sloppy          % Justification moins stricte : des mots ne dépasseront pas des paragraphes
  %\cleardoublepage
  \begin{onehalfspace}
  \pagestyle{fancy}
  
  \chapter{Introduction}
\thispagestyle{fancy}
Le but de ce projet est de nous familiariser avec les environnement 3D et de faire une première approche des système multi-agent.

 



  
  
  \chapter{Conclusion}


  \begin{spacing}{2.0}
    \include{Bibliographie}
  \end{spacing}
  \chapter*{Annexe}



\end{onehalfspace}
\end{document}
